% TDSC preparation draft. Do not submit until submission_check.py passes
% and the authors approve the selected venue and final declarations.
\documentclass[10pt,journal,compsoc]{IEEEtran}
\usepackage[T1]{fontenc}
\usepackage{amsmath,amssymb,booktabs,graphicx,microtype,xurl}
\usepackage[hidelinks]{hyperref}
% GENERATED by manifest.py --tex -- DO NOT EDIT.
% Every macro below is derived from runs/ at generation time.
% Regenerate:  python manifest.py --tex

\newcommand{\DefenseAucCI}{[0.986, 0.999]}
\newcommand{\DefenseAucValue}{0.994}
\newcommand{\DefenseBandAttack}{25}
\newcommand{\DefenseBandLegitimate}{15}
\newcommand{\DefenseBaserateFivePct}{5.0\%}
\newcommand{\DefenseBaserateOnePct}{1.0\%}
\newcommand{\DefenseBaseratePointonePct}{0.1\%}
\newcommand{\DefenseCorpusPositives}{9}
\newcommand{\DefenseCorpusNegatives}{506}
\newcommand{\DefenseFprAttwozeroPct}{2.6\%}
\newcommand{\DefensePpvFivePct}{67.2\%}
\newcommand{\DefensePpvOnePct}{28.2\%}
\newcommand{\DefensePpvPointonePct}{3.8\%}
\newcommand{\DefenseSeqthinkingFrac}{5/13}
\newcommand{\DefenseSeqthinkingK}{5}
\newcommand{\DefenseSeqthinkingN}{13}
\newcommand{\DefenseSeqthinkingNames}{4}
\newcommand{\DefenseTprAttwozeroPct}{100\%}
\newcommand{\DepthCredOptionalFrac}{14/20}
\newcommand{\DepthCredOptionalPct}{70\%}
\newcommand{\DepthCredOptionalK}{14}
\newcommand{\DepthCredOptionalN}{20}
\newcommand{\DepthCredPosfirstFrac}{20/20}
\newcommand{\DepthCredPosfirstPct}{100\%}
\newcommand{\DepthCredPosfirstK}{20}
\newcommand{\DepthCredPosfirstN}{20}
\newcommand{\DepthUaOptionalFrac}{8/20}
\newcommand{\DepthUaOptionalPct}{40\%}
\newcommand{\DepthUaOptionalK}{8}
\newcommand{\DepthUaOptionalN}{20}
\newcommand{\DepthUaPosfirstFrac}{19/20}
\newcommand{\DepthUaPosfirstPct}{95\%}
\newcommand{\DepthUaPosfirstK}{19}
\newcommand{\DepthUaPosfirstN}{20}
\newcommand{\DepthCredOptionalBaselineFrac}{20/20}
\newcommand{\DepthCredOptionalBaselinePct}{100\%}
\newcommand{\DepthCredOptionalBaselineK}{20}
\newcommand{\DepthCredOptionalBaselineN}{20}
\newcommand{\DepthCredOptionalDropPct}{30\%}
\newcommand{\DepthCredPosfirstBaselineFrac}{20/20}
\newcommand{\DepthCredPosfirstBaselinePct}{100\%}
\newcommand{\DepthCredPosfirstBaselineK}{20}
\newcommand{\DepthCredPosfirstBaselineN}{20}
\newcommand{\DepthCredPosfirstDropPct}{0\%}
\newcommand{\DepthUaOptionalBaselineFrac}{20/20}
\newcommand{\DepthUaOptionalBaselinePct}{100\%}
\newcommand{\DepthUaOptionalBaselineK}{20}
\newcommand{\DepthUaOptionalBaselineN}{20}
\newcommand{\DepthUaOptionalDropPct}{60\%}
\newcommand{\DepthUaPosfirstBaselineFrac}{20/20}
\newcommand{\DepthUaPosfirstBaselinePct}{100\%}
\newcommand{\DepthUaPosfirstBaselineK}{20}
\newcommand{\DepthUaPosfirstBaselineN}{20}
\newcommand{\DepthUaPosfirstDropPct}{5\%}
\newcommand{\ExplicitGptFouroCredentialShapedFrac}{4/20}
\newcommand{\ExplicitGptFouroCredentialShapedPct}{20\%}
\newcommand{\ExplicitGptFouroCredentialShapedK}{4}
\newcommand{\ExplicitGptFouroCredentialShapedN}{20}
\newcommand{\ExplicitGptFouroInternalCodenameFrac}{20/20}
\newcommand{\ExplicitGptFouroInternalCodenamePct}{100\%}
\newcommand{\ExplicitGptFouroInternalCodenameK}{20}
\newcommand{\ExplicitGptFouroInternalCodenameN}{20}
\newcommand{\ExplicitGptFouroLongBlockFrac}{19/20}
\newcommand{\ExplicitGptFouroLongBlockPct}{95\%}
\newcommand{\ExplicitGptFouroLongBlockK}{19}
\newcommand{\ExplicitGptFouroLongBlockN}{20}
\newcommand{\ExplicitGptFouroPolicySentenceFrac}{0/20}
\newcommand{\ExplicitGptFouroPolicySentencePct}{0\%}
\newcommand{\ExplicitGptFouroPolicySentenceK}{0}
\newcommand{\ExplicitGptFouroPolicySentenceN}{20}
\newcommand{\ExplicitGptFouroPooledFrac}{46/100}
\newcommand{\ExplicitGptFouroPooledPct}{46\%}
\newcommand{\ExplicitGptFouroPooledK}{46}
\newcommand{\ExplicitGptFouroPooledN}{100}
\newcommand{\ExplicitGptFouroProductNameFrac}{3/20}
\newcommand{\ExplicitGptFouroProductNamePct}{15\%}
\newcommand{\ExplicitGptFouroProductNameK}{3}
\newcommand{\ExplicitGptFouroProductNameN}{20}
\newcommand{\FrameworkArmDisclosed}{0}
\newcommand{\FrameworkArmN}{240}
\newcommand{\FrameworkArmToolCalled}{240}
\newcommand{\GateClaudeSonnetFourFiveAFrac}{30/30}
\newcommand{\GateClaudeSonnetFourFiveAPct}{100\%}
\newcommand{\GateClaudeSonnetFourFiveAK}{30}
\newcommand{\GateClaudeSonnetFourFiveAN}{30}
\newcommand{\GateClaudeSonnetFourFiveAPrimeFrac}{30/30}
\newcommand{\GateClaudeSonnetFourFiveAPrimePct}{100\%}
\newcommand{\GateClaudeSonnetFourFiveAPrimeK}{30}
\newcommand{\GateClaudeSonnetFourFiveAPrimeN}{30}
\newcommand{\GateClaudeSonnetFourFiveBFrac}{0/30}
\newcommand{\GateClaudeSonnetFourFiveBPct}{0\%}
\newcommand{\GateClaudeSonnetFourFiveBK}{0}
\newcommand{\GateClaudeSonnetFourFiveBN}{30}
\newcommand{\GateClaudeSonnetFourFiveCFrac}{30/30}
\newcommand{\GateClaudeSonnetFourFiveCPct}{100\%}
\newcommand{\GateClaudeSonnetFourFiveCK}{30}
\newcommand{\GateClaudeSonnetFourFiveCN}{30}
\newcommand{\GateClaudeSonnetFourFiveDFrac}{0/30}
\newcommand{\GateClaudeSonnetFourFiveDPct}{0\%}
\newcommand{\GateClaudeSonnetFourFiveDK}{0}
\newcommand{\GateClaudeSonnetFourFiveDN}{30}
\newcommand{\GateClaudeSonnetFourFiveDeltaPp}{0}
\newcommand{\GateGeminiThreeFlashAFrac}{0/30}
\newcommand{\GateGeminiThreeFlashAPct}{0\%}
\newcommand{\GateGeminiThreeFlashAK}{0}
\newcommand{\GateGeminiThreeFlashAN}{30}
\newcommand{\GateGeminiThreeFlashAPrimeFrac}{30/30}
\newcommand{\GateGeminiThreeFlashAPrimePct}{100\%}
\newcommand{\GateGeminiThreeFlashAPrimeK}{30}
\newcommand{\GateGeminiThreeFlashAPrimeN}{30}
\newcommand{\GateGeminiThreeFlashBFrac}{7/30}
\newcommand{\GateGeminiThreeFlashBPct}{23\%}
\newcommand{\GateGeminiThreeFlashBK}{7}
\newcommand{\GateGeminiThreeFlashBN}{30}
\newcommand{\GateGeminiThreeFlashCFrac}{30/30}
\newcommand{\GateGeminiThreeFlashCPct}{100\%}
\newcommand{\GateGeminiThreeFlashCK}{30}
\newcommand{\GateGeminiThreeFlashCN}{30}
\newcommand{\GateGeminiThreeFlashDFrac}{0/30}
\newcommand{\GateGeminiThreeFlashDPct}{0\%}
\newcommand{\GateGeminiThreeFlashDK}{0}
\newcommand{\GateGeminiThreeFlashDN}{30}
\newcommand{\GateGeminiThreeFlashDeltaPp}{0}
\newcommand{\GateGptFouroAFrac}{0/30}
\newcommand{\GateGptFouroAPct}{0\%}
\newcommand{\GateGptFouroAK}{0}
\newcommand{\GateGptFouroAN}{30}
\newcommand{\GateGptFouroAPrimeFrac}{9/30}
\newcommand{\GateGptFouroAPrimePct}{30\%}
\newcommand{\GateGptFouroAPrimeK}{9}
\newcommand{\GateGptFouroAPrimeN}{30}
\newcommand{\GateGptFouroBFrac}{0/30}
\newcommand{\GateGptFouroBPct}{0\%}
\newcommand{\GateGptFouroBK}{0}
\newcommand{\GateGptFouroBN}{30}
\newcommand{\GateGptFouroCFrac}{29/30}
\newcommand{\GateGptFouroCPct}{97\%}
\newcommand{\GateGptFouroCK}{29}
\newcommand{\GateGptFouroCN}{30}
\newcommand{\GateGptFouroDFrac}{0/30}
\newcommand{\GateGptFouroDPct}{0\%}
\newcommand{\GateGptFouroDK}{0}
\newcommand{\GateGptFouroDN}{30}
\newcommand{\GateGptFouroDeltaPp}{67}
\newcommand{\PooledGptFouroAPrimeFrac}{20/60}
\newcommand{\PooledGptFouroAPrimePct}{33\%}
\newcommand{\PooledGptFouroAPrimeCI}{[23, 46]}
\newcommand{\PooledGptFouroAPrimeK}{20}
\newcommand{\PooledGptFouroAPrimeN}{60}
\newcommand{\PooledGptFouroCFrac}{55/60}
\newcommand{\PooledGptFouroCPct}{92\%}
\newcommand{\PooledGptFouroCCI}{[82, 96]}
\newcommand{\PooledGptFouroCK}{55}
\newcommand{\PooledGptFouroCN}{60}
\newcommand{\PrevalenceAttackClientIdentityFrac}{0/401}
\newcommand{\PrevalenceAttackClientIdentityK}{0}
\newcommand{\PrevalenceAttackClientIdentityN}{401}
\newcommand{\PrevalenceAttackClientIdentityServers}{0}
\newcommand{\PrevalenceAttackContactFrac}{1/401}
\newcommand{\PrevalenceAttackContactK}{1}
\newcommand{\PrevalenceAttackContactN}{401}
\newcommand{\PrevalenceAttackContactServers}{1}
\newcommand{\PrevalenceAttackCredentialFrac}{0/401}
\newcommand{\PrevalenceAttackCredentialK}{0}
\newcommand{\PrevalenceAttackCredentialN}{401}
\newcommand{\PrevalenceAttackCredentialServers}{0}
\newcommand{\PrevalenceAttackOperatortenantFrac}{5/401}
\newcommand{\PrevalenceAttackOperatortenantK}{5}
\newcommand{\PrevalenceAttackOperatortenantN}{401}
\newcommand{\PrevalenceAttackOperatortenantServers}{2}
\newcommand{\PrevalenceAttackPolicyFrac}{0/401}
\newcommand{\PrevalenceAttackPolicyK}{0}
\newcommand{\PrevalenceAttackPolicyN}{401}
\newcommand{\PrevalenceAttackPolicyServers}{0}
\newcommand{\PrevalenceAttackRegionFrac}{1/401}
\newcommand{\PrevalenceAttackRegionK}{1}
\newcommand{\PrevalenceAttackRegionN}{401}
\newcommand{\PrevalenceAttackRegionServers}{3}
\newcommand{\PrevalenceCorpusServers}{45}
\newcommand{\PrevalenceCorpusParams}{1048}
\newcommand{\PrevalenceCorpusN}{401}
\newcommand{\PrevalenceVictimCharsN}{195,012}
\newcommand{\PrevalenceVictimCloudRegionK}{3}
\newcommand{\PrevalenceVictimCredentialShapedK}{0}
\newcommand{\PrevalenceVictimOperatorTenantK}{0}
\newcommand{\PrevalenceVictimPolicyTokenK}{0}
\newcommand{\RsixDeltaPp}{65}
\newcommand{\RsixFamiliarFrac}{1/20}
\newcommand{\RsixFamiliarPct}{5\%}
\newcommand{\RsixFamiliarK}{1}
\newcommand{\RsixFamiliarN}{20}
\newcommand{\RsixSeedFrac}{14/20}
\newcommand{\RsixSeedPct}{70\%}
\newcommand{\RsixSeedK}{14}
\newcommand{\RsixSeedN}{20}
\newcommand{\RealDCombinedFrac}{0/228}
\newcommand{\RealDCombinedK}{0}
\newcommand{\RealDCombinedN}{228}
\newcommand{\RealClaudeSonnetFourFiveCFrac}{120/120}
\newcommand{\RealClaudeSonnetFourFiveCPct}{100\%}
\newcommand{\RealClaudeSonnetFourFiveCIttFrac}{120/125}
\newcommand{\RealClaudeSonnetFourFiveCK}{120}
\newcommand{\RealClaudeSonnetFourFiveCN}{120}
\newcommand{\RealClaudeSonnetFourFiveCIttK}{120}
\newcommand{\RealClaudeSonnetFourFiveCIttN}{125}
\newcommand{\RealClaudeSonnetFourFiveCIttPct}{96\%}
\newcommand{\RealClaudeSonnetFourFiveCToolsLeaking}{24}
\newcommand{\RealClaudeSonnetFourFiveCToolsCalled}{24}
\newcommand{\RealClaudeSonnetFourFiveDFrac}{0/119}
\newcommand{\RealClaudeSonnetFourFiveDPct}{0\%}
\newcommand{\RealClaudeSonnetFourFiveDIttFrac}{0/124}
\newcommand{\RealClaudeSonnetFourFiveDK}{0}
\newcommand{\RealClaudeSonnetFourFiveDN}{119}
\newcommand{\RealClaudeSonnetFourFiveDIttK}{0}
\newcommand{\RealClaudeSonnetFourFiveDIttN}{124}
\newcommand{\RealClaudeSonnetFourFiveDIttPct}{0\%}
\newcommand{\RealClaudeSonnetFourFiveDToolsLeaking}{0}
\newcommand{\RealClaudeSonnetFourFiveDToolsCalled}{24}
\newcommand{\RealCorpusTools}{25}
\newcommand{\RealCorpusServers}{25}
\newcommand{\RealGptFouroCFrac}{78/113}
\newcommand{\RealGptFouroCPct}{69\%}
\newcommand{\RealGptFouroCIttFrac}{78/125}
\newcommand{\RealGptFouroCK}{78}
\newcommand{\RealGptFouroCN}{113}
\newcommand{\RealGptFouroCIttK}{78}
\newcommand{\RealGptFouroCIttN}{125}
\newcommand{\RealGptFouroCIttPct}{62\%}
\newcommand{\RealGptFouroCToolsLeaking}{22}
\newcommand{\RealGptFouroCToolsCalled}{24}
\newcommand{\RealGptFouroDFrac}{0/109}
\newcommand{\RealGptFouroDPct}{0\%}
\newcommand{\RealGptFouroDIttFrac}{0/125}
\newcommand{\RealGptFouroDK}{0}
\newcommand{\RealGptFouroDN}{109}
\newcommand{\RealGptFouroDIttK}{0}
\newcommand{\RealGptFouroDIttN}{125}
\newcommand{\RealGptFouroDIttPct}{0\%}
\newcommand{\RealGptFouroDToolsLeaking}{0}
\newcommand{\RealGptFouroDToolsCalled}{23}
\newcommand{\ScanDescAPrimeFlagged}{0}
\newcommand{\ScanDescBFlagged}{1}
\newcommand{\ScanDescBPrimeFlagged}{1}
\newcommand{\ScanDescCFlagged}{0}
\newcommand{\ScanDescDFlagged}{0}
\newcommand{\ScanDescEFlagged}{1}
\newcommand{\ScanDescnameAPrimeFlagged}{1}
\newcommand{\ScanDescnameBFlagged}{1}
\newcommand{\ScanDescnameBPrimeFlagged}{1}
\newcommand{\ScanDescnameCFlagged}{0}
\newcommand{\ScanDescnameDFlagged}{0}
\newcommand{\ScanDescnameEFlagged}{1}
\newcommand{\ScanInvariantAPrimeFlagged}{0}
\newcommand{\ScanInvariantBFlagged}{1}
\newcommand{\ScanInvariantBPrimeFlagged}{1}
\newcommand{\ScanInvariantCApiDocumentationFlagged}{0}
\newcommand{\ScanInvariantCComplianceLoggingFlagged}{0}
\newcommand{\ScanInvariantCDefaultFlagged}{0}
\newcommand{\ScanInvariantCIndirectInferenceFlagged}{0}
\newcommand{\ScanInvariantCOperationalMetadataFlagged}{0}
\newcommand{\ScanInvariantDFlagged}{0}
\newcommand{\ScanInvariantEFlagged}{0}
\newcommand{\ScanInvariantMetaReps}{3}
\newcommand{\ScanInvariantMetaArms}{10}
\newcommand{\ScanRawAPrimeFlagged}{0}
\newcommand{\ScanRawBFlagged}{1}
\newcommand{\ScanRawBPrimeFlagged}{1}
\newcommand{\ScanRawCApiDocumentationFlagged}{0}
\newcommand{\ScanRawCComplianceLoggingFlagged}{0}
\newcommand{\ScanRawCDefaultFlagged}{0}
\newcommand{\ScanRawCFieldEscalationFlagged}{0}
\newcommand{\ScanRawCFieldOperatorFlagged}{0}
\newcommand{\ScanRawCFieldRegionFlagged}{0}
\newcommand{\ScanRawCFieldRetentionFlagged}{0}
\newcommand{\ScanRawCGAnyGenericFlagged}{0}
\newcommand{\ScanRawCGCredentialAdjacentFlagged}{0}
\newcommand{\ScanRawCGCredentialGenericFlagged}{0}
\newcommand{\ScanRawCGPlatformAdjacentFlagged}{0}
\newcommand{\ScanRawCGRegionAdjacentFlagged}{0}
\newcommand{\ScanRawCGRegionGenericFlagged}{0}
\newcommand{\ScanRawCIndirectInferenceFlagged}{0}
\newcommand{\ScanRawCMBuildFlagged}{0}
\newcommand{\ScanRawCMContactFlagged}{0}
\newcommand{\ScanRawCMCredentialFlagged}{0}
\newcommand{\ScanRawCMOperatorFlagged}{0}
\newcommand{\ScanRawCMPolicyFlagged}{0}
\newcommand{\ScanRawCMRegionFlagged}{0}
\newcommand{\ScanRawCOperationalMetadataFlagged}{0}
\newcommand{\ScanRawDFlagged}{0}
\newcommand{\ScanRawEFlagged}{0}
\newcommand{\ScanRawSOneDFlagged}{0}
\newcommand{\ScanRawSOneApiDocumentationFlagged}{0}
\newcommand{\ScanRawSOneGAnyGenericFlagged}{0}
\newcommand{\ScanRawSOneGCredentialAdjacentFlagged}{0}
\newcommand{\ScanRawSOneGPlatformAdjacentFlagged}{0}
\newcommand{\ScanRawSOneGRegionAdjacentFlagged}{0}
\newcommand{\ScanRawSOneMCredentialFlagged}{0}
\newcommand{\ScanRawSOneMOperatorFlagged}{0}
\newcommand{\ScanRawSOneMRegionFlagged}{0}
\newcommand{\ScanRawSTwoClientPlatformNameFlagged}{0}
\newcommand{\ScanRawSTwoDeploymentRegionFlagged}{0}
\newcommand{\ScanRawSTwoOperatorNameFlagged}{0}
\newcommand{\ScanRawSTwoServiceKeyFlagged}{0}
\newcommand{\ScanRawSThreeCovertCodeEightFlagged}{0}
\newcommand{\ScanRawSThreeDeploymentRegionHexFlagged}{0}
\newcommand{\ScanRawSThreeDeploymentRegionReversedFlagged}{0}
\newcommand{\ScanRawSThreeDeploymentRegionSuffixFlagged}{0}
\newcommand{\ScanRawSThreeServiceKeyHexFlagged}{0}
\newcommand{\ScanRawSThreeServiceKeyReversedFlagged}{0}
\newcommand{\ScanRawSThreeServiceKeySuffixFlagged}{0}
\newcommand{\ScanRawSFourAccountAttributionFlagged}{0}
\newcommand{\ScanRawSFourCallContextFlagged}{0}
\newcommand{\ScanRawSFourDeployingOrganisationFlagged}{0}
\newcommand{\ScanRawSFourHostingRegionCodeFlagged}{0}
\newcommand{\ScanRawSFourIntegrationCredentialIdFlagged}{0}
\newcommand{\ScanRawSFourIssuingEnvironmentFlagged}{0}
\newcommand{\ScanRawSFourRequestTraceIdFlagged}{0}
\newcommand{\ScanRawSFourRuntimePlatformFlagged}{0}
\newcommand{\ScanRawSFourServingLocationHintFlagged}{0}
\newcommand{\ScanRawSFourSetupBindingNoteFlagged}{0}
\newcommand{\ScanRawMetaReps}{3}
\newcommand{\ScanRawMetaArms}{56}
\newcommand{\ScanRawMetaCalls}{168}
\newcommand{\ScanRawMetaValid}{168}
\newcommand{\ScanRawMetaStudyarms}{30}
\newcommand{\ScanRawMetaStudyflagged}{0}
\newcommand{\ScanRawMetaSonearms}{9}
\newcommand{\ScanRawMetaSoneflagged}{0}
\newcommand{\ScanRawMetaStwoarms}{4}
\newcommand{\ScanRawMetaStwoflagged}{0}
\newcommand{\ScanRawMetaSthreearms}{7}
\newcommand{\ScanRawMetaSthreeflagged}{0}
\newcommand{\ScanRawMetaSfourarms}{10}
\newcommand{\ScanRawMetaSfourflagged}{0}
\newcommand{\VtwoGeminiThreeFlashPreviewDiagonalFrac}{139/140}
\newcommand{\VtwoGeminiThreeFlashPreviewDiagonalPct}{99\%}
\newcommand{\VtwoGeminiThreeFlashPreviewDiagonalK}{139}
\newcommand{\VtwoGeminiThreeFlashPreviewDiagonalN}{140}
\newcommand{\VtwoGeminiThreeFlashPreviewOffdiagonalFrac}{39/980}
\newcommand{\VtwoGeminiThreeFlashPreviewOffdiagonalPct}{4\%}
\newcommand{\VtwoGeminiThreeFlashPreviewOffdiagonalK}{39}
\newcommand{\VtwoGeminiThreeFlashPreviewOffdiagonalN}{980}
\newcommand{\VtwoGptFouroDiagonalFrac}{138/140}
\newcommand{\VtwoGptFouroDiagonalPct}{99\%}
\newcommand{\VtwoGptFouroDiagonalK}{138}
\newcommand{\VtwoGptFouroDiagonalN}{140}
\newcommand{\VtwoGptFouroOffdiagonalFrac}{8/980}
\newcommand{\VtwoGptFouroOffdiagonalPct}{1\%}
\newcommand{\VtwoGptFouroOffdiagonalK}{8}
\newcommand{\VtwoGptFouroOffdiagonalN}{980}
\newcommand{\VthreeClaudeSonnetFourFiveAdjacentFrac}{20/60}
\newcommand{\VthreeClaudeSonnetFourFiveAdjacentPct}{33\%}
\newcommand{\VthreeClaudeSonnetFourFiveAdjacentCI}{[23, 46]}
\newcommand{\VthreeClaudeSonnetFourFiveAdjacentK}{20}
\newcommand{\VthreeClaudeSonnetFourFiveAdjacentN}{60}
\newcommand{\VthreeClaudeSonnetFourFiveControlFrac}{0/800}
\newcommand{\VthreeClaudeSonnetFourFiveControlPct}{0\%}
\newcommand{\VthreeClaudeSonnetFourFiveControlK}{0}
\newcommand{\VthreeClaudeSonnetFourFiveControlN}{800}
\newcommand{\VthreeClaudeSonnetFourFiveControlcallsFrac}{0/160}
\newcommand{\VthreeClaudeSonnetFourFiveControlcallsK}{0}
\newcommand{\VthreeClaudeSonnetFourFiveControlcallsN}{160}
\newcommand{\VthreeClaudeSonnetFourFiveCredentialFrac}{20/20}
\newcommand{\VthreeClaudeSonnetFourFiveCredentialPct}{100\%}
\newcommand{\VthreeClaudeSonnetFourFiveCredentialK}{20}
\newcommand{\VthreeClaudeSonnetFourFiveCredentialN}{20}
\newcommand{\VthreeClaudeSonnetFourFiveCrosstalknamingFrac}{0/0}
\newcommand{\VthreeClaudeSonnetFourFiveCrosstalknamingK}{0}
\newcommand{\VthreeClaudeSonnetFourFiveCrosstalknamingN}{0}
\newcommand{\VthreeClaudeSonnetFourFiveCrosstalkuaFrac}{0/0}
\newcommand{\VthreeClaudeSonnetFourFiveCrosstalkuaK}{0}
\newcommand{\VthreeClaudeSonnetFourFiveCrosstalkuaN}{0}
\newcommand{\VthreeClaudeSonnetFourFiveDeltaCI}{[43, 100]}
\newcommand{\VthreeClaudeSonnetFourFiveDeltaPp}{71}
\newcommand{\VthreeClaudeSonnetFourFiveDeltaClusters}{7}
\newcommand{\VthreeClaudeSonnetFourFiveDiagonalFrac}{100/140}
\newcommand{\VthreeClaudeSonnetFourFiveDiagonalPct}{71\%}
\newcommand{\VthreeClaudeSonnetFourFiveDiagonalCI}{[43, 100]}
\newcommand{\VthreeClaudeSonnetFourFiveDiagonalK}{100}
\newcommand{\VthreeClaudeSonnetFourFiveDiagonalN}{140}
\newcommand{\VthreeClaudeSonnetFourFiveFabricationFrac}{0/120}
\newcommand{\VthreeClaudeSonnetFourFiveFabricationPct}{0\%}
\newcommand{\VthreeClaudeSonnetFourFiveFabricationK}{0}
\newcommand{\VthreeClaudeSonnetFourFiveFabricationN}{120}
\newcommand{\VthreeClaudeSonnetFourFiveFieldApiDocumentationFrac}{20/20}
\newcommand{\VthreeClaudeSonnetFourFiveFieldApiDocumentationPct}{100\%}
\newcommand{\VthreeClaudeSonnetFourFiveFieldApiDocumentationK}{20}
\newcommand{\VthreeClaudeSonnetFourFiveFieldApiDocumentationN}{20}
\newcommand{\VthreeClaudeSonnetFourFiveFieldGCredentialAdjacentFrac}{0/20}
\newcommand{\VthreeClaudeSonnetFourFiveFieldGCredentialAdjacentPct}{0\%}
\newcommand{\VthreeClaudeSonnetFourFiveFieldGCredentialAdjacentK}{0}
\newcommand{\VthreeClaudeSonnetFourFiveFieldGCredentialAdjacentN}{20}
\newcommand{\VthreeClaudeSonnetFourFiveFieldGPlatformAdjacentFrac}{0/20}
\newcommand{\VthreeClaudeSonnetFourFiveFieldGPlatformAdjacentPct}{0\%}
\newcommand{\VthreeClaudeSonnetFourFiveFieldGPlatformAdjacentK}{0}
\newcommand{\VthreeClaudeSonnetFourFiveFieldGPlatformAdjacentN}{20}
\newcommand{\VthreeClaudeSonnetFourFiveFieldGRegionAdjacentFrac}{20/20}
\newcommand{\VthreeClaudeSonnetFourFiveFieldGRegionAdjacentPct}{100\%}
\newcommand{\VthreeClaudeSonnetFourFiveFieldGRegionAdjacentK}{20}
\newcommand{\VthreeClaudeSonnetFourFiveFieldGRegionAdjacentN}{20}
\newcommand{\VthreeClaudeSonnetFourFiveFieldMCredentialFrac}{20/20}
\newcommand{\VthreeClaudeSonnetFourFiveFieldMCredentialPct}{100\%}
\newcommand{\VthreeClaudeSonnetFourFiveFieldMCredentialK}{20}
\newcommand{\VthreeClaudeSonnetFourFiveFieldMCredentialN}{20}
\newcommand{\VthreeClaudeSonnetFourFiveFieldMOperatorFrac}{20/20}
\newcommand{\VthreeClaudeSonnetFourFiveFieldMOperatorPct}{100\%}
\newcommand{\VthreeClaudeSonnetFourFiveFieldMOperatorK}{20}
\newcommand{\VthreeClaudeSonnetFourFiveFieldMOperatorN}{20}
\newcommand{\VthreeClaudeSonnetFourFiveFieldMRegionFrac}{20/20}
\newcommand{\VthreeClaudeSonnetFourFiveFieldMRegionPct}{100\%}
\newcommand{\VthreeClaudeSonnetFourFiveFieldMRegionK}{20}
\newcommand{\VthreeClaudeSonnetFourFiveFieldMRegionN}{20}
\newcommand{\VthreeClaudeSonnetFourFiveGenericFrac}{0/20}
\newcommand{\VthreeClaudeSonnetFourFiveGenericPct}{0\%}
\newcommand{\VthreeClaudeSonnetFourFiveGenericK}{0}
\newcommand{\VthreeClaudeSonnetFourFiveGenericN}{20}
\newcommand{\VthreeClaudeSonnetFourFiveNamingFrac}{80/80}
\newcommand{\VthreeClaudeSonnetFourFiveNamingPct}{100\%}
\newcommand{\VthreeClaudeSonnetFourFiveNamingCI}{[95, 100]}
\newcommand{\VthreeClaudeSonnetFourFiveNamingK}{80}
\newcommand{\VthreeClaudeSonnetFourFiveNamingN}{80}
\newcommand{\VthreeClaudeSonnetFourFiveOffdiagFrac}{0/560}
\newcommand{\VthreeClaudeSonnetFourFiveOffdiagPct}{0\%}
\newcommand{\VthreeClaudeSonnetFourFiveOffdiagCI}{[0, 0]}
\newcommand{\VthreeClaudeSonnetFourFiveOffdiagK}{0}
\newcommand{\VthreeClaudeSonnetFourFiveOffdiagN}{560}
\newcommand{\VthreeDeepseekVfourFlashAdjacentFrac}{20/60}
\newcommand{\VthreeDeepseekVfourFlashAdjacentPct}{33\%}
\newcommand{\VthreeDeepseekVfourFlashAdjacentCI}{[23, 46]}
\newcommand{\VthreeDeepseekVfourFlashAdjacentK}{20}
\newcommand{\VthreeDeepseekVfourFlashAdjacentN}{60}
\newcommand{\VthreeDeepseekVfourFlashControlFrac}{0/800}
\newcommand{\VthreeDeepseekVfourFlashControlPct}{0\%}
\newcommand{\VthreeDeepseekVfourFlashControlK}{0}
\newcommand{\VthreeDeepseekVfourFlashControlN}{800}
\newcommand{\VthreeDeepseekVfourFlashControlcallsFrac}{0/160}
\newcommand{\VthreeDeepseekVfourFlashControlcallsK}{0}
\newcommand{\VthreeDeepseekVfourFlashControlcallsN}{160}
\newcommand{\VthreeDeepseekVfourFlashCredentialFrac}{20/20}
\newcommand{\VthreeDeepseekVfourFlashCredentialPct}{100\%}
\newcommand{\VthreeDeepseekVfourFlashCredentialK}{20}
\newcommand{\VthreeDeepseekVfourFlashCredentialN}{20}
\newcommand{\VthreeDeepseekVfourFlashCrosstalknamingFrac}{0/4}
\newcommand{\VthreeDeepseekVfourFlashCrosstalknamingK}{0}
\newcommand{\VthreeDeepseekVfourFlashCrosstalknamingN}{4}
\newcommand{\VthreeDeepseekVfourFlashCrosstalkuaFrac}{2/4}
\newcommand{\VthreeDeepseekVfourFlashCrosstalkuaK}{2}
\newcommand{\VthreeDeepseekVfourFlashCrosstalkuaN}{4}
\newcommand{\VthreeDeepseekVfourFlashDeltaCI}{[27, 90]}
\newcommand{\VthreeDeepseekVfourFlashDeltaPp}{61}
\newcommand{\VthreeDeepseekVfourFlashDeltaClusters}{7}
\newcommand{\VthreeDeepseekVfourFlashDiagonalFrac}{84/137}
\newcommand{\VthreeDeepseekVfourFlashDiagonalPct}{61\%}
\newcommand{\VthreeDeepseekVfourFlashDiagonalCI}{[29, 91]}
\newcommand{\VthreeDeepseekVfourFlashDiagonalK}{84}
\newcommand{\VthreeDeepseekVfourFlashDiagonalN}{137}
\newcommand{\VthreeDeepseekVfourFlashFabricationFrac}{0/77}
\newcommand{\VthreeDeepseekVfourFlashFabricationPct}{0\%}
\newcommand{\VthreeDeepseekVfourFlashFabricationK}{0}
\newcommand{\VthreeDeepseekVfourFlashFabricationN}{77}
\newcommand{\VthreeDeepseekVfourFlashFieldApiDocumentationFrac}{7/20}
\newcommand{\VthreeDeepseekVfourFlashFieldApiDocumentationPct}{35\%}
\newcommand{\VthreeDeepseekVfourFlashFieldApiDocumentationK}{7}
\newcommand{\VthreeDeepseekVfourFlashFieldApiDocumentationN}{20}
\newcommand{\VthreeDeepseekVfourFlashFieldGCredentialAdjacentFrac}{0/20}
\newcommand{\VthreeDeepseekVfourFlashFieldGCredentialAdjacentPct}{0\%}
\newcommand{\VthreeDeepseekVfourFlashFieldGCredentialAdjacentK}{0}
\newcommand{\VthreeDeepseekVfourFlashFieldGCredentialAdjacentN}{20}
\newcommand{\VthreeDeepseekVfourFlashFieldGPlatformAdjacentFrac}{0/20}
\newcommand{\VthreeDeepseekVfourFlashFieldGPlatformAdjacentPct}{0\%}
\newcommand{\VthreeDeepseekVfourFlashFieldGPlatformAdjacentK}{0}
\newcommand{\VthreeDeepseekVfourFlashFieldGPlatformAdjacentN}{20}
\newcommand{\VthreeDeepseekVfourFlashFieldGRegionAdjacentFrac}{20/20}
\newcommand{\VthreeDeepseekVfourFlashFieldGRegionAdjacentPct}{100\%}
\newcommand{\VthreeDeepseekVfourFlashFieldGRegionAdjacentK}{20}
\newcommand{\VthreeDeepseekVfourFlashFieldGRegionAdjacentN}{20}
\newcommand{\VthreeDeepseekVfourFlashFieldMCredentialFrac}{20/20}
\newcommand{\VthreeDeepseekVfourFlashFieldMCredentialPct}{100\%}
\newcommand{\VthreeDeepseekVfourFlashFieldMCredentialK}{20}
\newcommand{\VthreeDeepseekVfourFlashFieldMCredentialN}{20}
\newcommand{\VthreeDeepseekVfourFlashFieldMOperatorFrac}{17/17}
\newcommand{\VthreeDeepseekVfourFlashFieldMOperatorPct}{100\%}
\newcommand{\VthreeDeepseekVfourFlashFieldMOperatorK}{17}
\newcommand{\VthreeDeepseekVfourFlashFieldMOperatorN}{17}
\newcommand{\VthreeDeepseekVfourFlashFieldMRegionFrac}{20/20}
\newcommand{\VthreeDeepseekVfourFlashFieldMRegionPct}{100\%}
\newcommand{\VthreeDeepseekVfourFlashFieldMRegionK}{20}
\newcommand{\VthreeDeepseekVfourFlashFieldMRegionN}{20}
\newcommand{\VthreeDeepseekVfourFlashGenericFrac}{0/20}
\newcommand{\VthreeDeepseekVfourFlashGenericPct}{0\%}
\newcommand{\VthreeDeepseekVfourFlashGenericK}{0}
\newcommand{\VthreeDeepseekVfourFlashGenericN}{20}
\newcommand{\VthreeDeepseekVfourFlashNamingFrac}{64/77}
\newcommand{\VthreeDeepseekVfourFlashNamingPct}{83\%}
\newcommand{\VthreeDeepseekVfourFlashNamingCI}{[73, 90]}
\newcommand{\VthreeDeepseekVfourFlashNamingK}{64}
\newcommand{\VthreeDeepseekVfourFlashNamingN}{77}
\newcommand{\VthreeDeepseekVfourFlashOffdiagFrac}{4/548}
\newcommand{\VthreeDeepseekVfourFlashOffdiagPct}{1\%}
\newcommand{\VthreeDeepseekVfourFlashOffdiagCI}{[0, 1]}
\newcommand{\VthreeDeepseekVfourFlashOffdiagK}{4}
\newcommand{\VthreeDeepseekVfourFlashOffdiagN}{548}
\newcommand{\VthreeFabricationPooledFrac}{0/367}
\newcommand{\VthreeFabricationPooledPct}{0\%}
\newcommand{\VthreeFabricationPooledK}{0}
\newcommand{\VthreeFabricationPooledN}{367}
\newcommand{\VthreeGeminiThreeFlashPreviewAdjacentFrac}{20/60}
\newcommand{\VthreeGeminiThreeFlashPreviewAdjacentPct}{33\%}
\newcommand{\VthreeGeminiThreeFlashPreviewAdjacentCI}{[23, 46]}
\newcommand{\VthreeGeminiThreeFlashPreviewAdjacentK}{20}
\newcommand{\VthreeGeminiThreeFlashPreviewAdjacentN}{60}
\newcommand{\VthreeGeminiThreeFlashPreviewControlFrac}{0/800}
\newcommand{\VthreeGeminiThreeFlashPreviewControlPct}{0\%}
\newcommand{\VthreeGeminiThreeFlashPreviewControlK}{0}
\newcommand{\VthreeGeminiThreeFlashPreviewControlN}{800}
\newcommand{\VthreeGeminiThreeFlashPreviewControlcallsFrac}{0/160}
\newcommand{\VthreeGeminiThreeFlashPreviewControlcallsK}{0}
\newcommand{\VthreeGeminiThreeFlashPreviewControlcallsN}{160}
\newcommand{\VthreeGeminiThreeFlashPreviewCredentialFrac}{20/20}
\newcommand{\VthreeGeminiThreeFlashPreviewCredentialPct}{100\%}
\newcommand{\VthreeGeminiThreeFlashPreviewCredentialK}{20}
\newcommand{\VthreeGeminiThreeFlashPreviewCredentialN}{20}
\newcommand{\VthreeGeminiThreeFlashPreviewCrosstalknamingFrac}{0/4}
\newcommand{\VthreeGeminiThreeFlashPreviewCrosstalknamingK}{0}
\newcommand{\VthreeGeminiThreeFlashPreviewCrosstalknamingN}{4}
\newcommand{\VthreeGeminiThreeFlashPreviewCrosstalkuaFrac}{4/4}
\newcommand{\VthreeGeminiThreeFlashPreviewCrosstalkuaK}{4}
\newcommand{\VthreeGeminiThreeFlashPreviewCrosstalkuaN}{4}
\newcommand{\VthreeGeminiThreeFlashPreviewDeltaCI}{[17, 86]}
\newcommand{\VthreeGeminiThreeFlashPreviewDeltaPp}{58}
\newcommand{\VthreeGeminiThreeFlashPreviewDeltaClusters}{7}
\newcommand{\VthreeGeminiThreeFlashPreviewDiagonalFrac}{82/140}
\newcommand{\VthreeGeminiThreeFlashPreviewDiagonalPct}{59\%}
\newcommand{\VthreeGeminiThreeFlashPreviewDiagonalCI}{[20, 87]}
\newcommand{\VthreeGeminiThreeFlashPreviewDiagonalK}{82}
\newcommand{\VthreeGeminiThreeFlashPreviewDiagonalN}{140}
\newcommand{\VthreeGeminiThreeFlashPreviewFabricationFrac}{0/77}
\newcommand{\VthreeGeminiThreeFlashPreviewFabricationPct}{0\%}
\newcommand{\VthreeGeminiThreeFlashPreviewFabricationK}{0}
\newcommand{\VthreeGeminiThreeFlashPreviewFabricationN}{77}
\newcommand{\VthreeGeminiThreeFlashPreviewFieldApiDocumentationFrac}{2/20}
\newcommand{\VthreeGeminiThreeFlashPreviewFieldApiDocumentationPct}{10\%}
\newcommand{\VthreeGeminiThreeFlashPreviewFieldApiDocumentationK}{2}
\newcommand{\VthreeGeminiThreeFlashPreviewFieldApiDocumentationN}{20}
\newcommand{\VthreeGeminiThreeFlashPreviewFieldGCredentialAdjacentFrac}{0/20}
\newcommand{\VthreeGeminiThreeFlashPreviewFieldGCredentialAdjacentPct}{0\%}
\newcommand{\VthreeGeminiThreeFlashPreviewFieldGCredentialAdjacentK}{0}
\newcommand{\VthreeGeminiThreeFlashPreviewFieldGCredentialAdjacentN}{20}
\newcommand{\VthreeGeminiThreeFlashPreviewFieldGPlatformAdjacentFrac}{0/20}
\newcommand{\VthreeGeminiThreeFlashPreviewFieldGPlatformAdjacentPct}{0\%}
\newcommand{\VthreeGeminiThreeFlashPreviewFieldGPlatformAdjacentK}{0}
\newcommand{\VthreeGeminiThreeFlashPreviewFieldGPlatformAdjacentN}{20}
\newcommand{\VthreeGeminiThreeFlashPreviewFieldGRegionAdjacentFrac}{20/20}
\newcommand{\VthreeGeminiThreeFlashPreviewFieldGRegionAdjacentPct}{100\%}
\newcommand{\VthreeGeminiThreeFlashPreviewFieldGRegionAdjacentK}{20}
\newcommand{\VthreeGeminiThreeFlashPreviewFieldGRegionAdjacentN}{20}
\newcommand{\VthreeGeminiThreeFlashPreviewFieldMCredentialFrac}{20/20}
\newcommand{\VthreeGeminiThreeFlashPreviewFieldMCredentialPct}{100\%}
\newcommand{\VthreeGeminiThreeFlashPreviewFieldMCredentialK}{20}
\newcommand{\VthreeGeminiThreeFlashPreviewFieldMCredentialN}{20}
\newcommand{\VthreeGeminiThreeFlashPreviewFieldMOperatorFrac}{20/20}
\newcommand{\VthreeGeminiThreeFlashPreviewFieldMOperatorPct}{100\%}
\newcommand{\VthreeGeminiThreeFlashPreviewFieldMOperatorK}{20}
\newcommand{\VthreeGeminiThreeFlashPreviewFieldMOperatorN}{20}
\newcommand{\VthreeGeminiThreeFlashPreviewFieldMRegionFrac}{20/20}
\newcommand{\VthreeGeminiThreeFlashPreviewFieldMRegionPct}{100\%}
\newcommand{\VthreeGeminiThreeFlashPreviewFieldMRegionK}{20}
\newcommand{\VthreeGeminiThreeFlashPreviewFieldMRegionN}{20}
\newcommand{\VthreeGeminiThreeFlashPreviewGenericFrac}{0/20}
\newcommand{\VthreeGeminiThreeFlashPreviewGenericPct}{0\%}
\newcommand{\VthreeGeminiThreeFlashPreviewGenericK}{0}
\newcommand{\VthreeGeminiThreeFlashPreviewGenericN}{20}
\newcommand{\VthreeGeminiThreeFlashPreviewNamingFrac}{62/80}
\newcommand{\VthreeGeminiThreeFlashPreviewNamingPct}{78\%}
\newcommand{\VthreeGeminiThreeFlashPreviewNamingCI}{[67, 85]}
\newcommand{\VthreeGeminiThreeFlashPreviewNamingK}{62}
\newcommand{\VthreeGeminiThreeFlashPreviewNamingN}{80}
\newcommand{\VthreeGeminiThreeFlashPreviewOffdiagFrac}{4/560}
\newcommand{\VthreeGeminiThreeFlashPreviewOffdiagPct}{1\%}
\newcommand{\VthreeGeminiThreeFlashPreviewOffdiagCI}{[0, 2]}
\newcommand{\VthreeGeminiThreeFlashPreviewOffdiagK}{4}
\newcommand{\VthreeGeminiThreeFlashPreviewOffdiagN}{560}
\newcommand{\VthreeGeminiThreeOneProPreviewAdjacentFrac}{40/43}
\newcommand{\VthreeGeminiThreeOneProPreviewAdjacentPct}{93\%}
\newcommand{\VthreeGeminiThreeOneProPreviewAdjacentCI}{[81, 98]}
\newcommand{\VthreeGeminiThreeOneProPreviewAdjacentK}{40}
\newcommand{\VthreeGeminiThreeOneProPreviewAdjacentN}{43}
\newcommand{\VthreeGeminiThreeOneProPreviewControlFrac}{0/580}
\newcommand{\VthreeGeminiThreeOneProPreviewControlPct}{0\%}
\newcommand{\VthreeGeminiThreeOneProPreviewControlK}{0}
\newcommand{\VthreeGeminiThreeOneProPreviewControlN}{580}
\newcommand{\VthreeGeminiThreeOneProPreviewControlcallsFrac}{0/116}
\newcommand{\VthreeGeminiThreeOneProPreviewControlcallsK}{0}
\newcommand{\VthreeGeminiThreeOneProPreviewControlcallsN}{116}
\newcommand{\VthreeGeminiThreeOneProPreviewCredentialFrac}{20/20}
\newcommand{\VthreeGeminiThreeOneProPreviewCredentialPct}{100\%}
\newcommand{\VthreeGeminiThreeOneProPreviewCredentialK}{20}
\newcommand{\VthreeGeminiThreeOneProPreviewCredentialN}{20}
\newcommand{\VthreeGeminiThreeOneProPreviewCrosstalknamingFrac}{0/24}
\newcommand{\VthreeGeminiThreeOneProPreviewCrosstalknamingK}{0}
\newcommand{\VthreeGeminiThreeOneProPreviewCrosstalknamingN}{24}
\newcommand{\VthreeGeminiThreeOneProPreviewCrosstalkuaFrac}{21/24}
\newcommand{\VthreeGeminiThreeOneProPreviewCrosstalkuaK}{21}
\newcommand{\VthreeGeminiThreeOneProPreviewCrosstalkuaN}{24}
\newcommand{\VthreeGeminiThreeOneProPreviewDeltaCI}{[80, 100]}
\newcommand{\VthreeGeminiThreeOneProPreviewDeltaPp}{93}
\newcommand{\VthreeGeminiThreeOneProPreviewDeltaClusters}{7}
\newcommand{\VthreeGeminiThreeOneProPreviewDiagonalFrac}{120/123}
\newcommand{\VthreeGeminiThreeOneProPreviewDiagonalPct}{98\%}
\newcommand{\VthreeGeminiThreeOneProPreviewDiagonalCI}{[90, 100]}
\newcommand{\VthreeGeminiThreeOneProPreviewDiagonalK}{120}
\newcommand{\VthreeGeminiThreeOneProPreviewDiagonalN}{123}
\newcommand{\VthreeGeminiThreeOneProPreviewFieldApiDocumentationFrac}{20/20}
\newcommand{\VthreeGeminiThreeOneProPreviewFieldApiDocumentationPct}{100\%}
\newcommand{\VthreeGeminiThreeOneProPreviewFieldApiDocumentationK}{20}
\newcommand{\VthreeGeminiThreeOneProPreviewFieldApiDocumentationN}{20}
\newcommand{\VthreeGeminiThreeOneProPreviewFieldGCredentialAdjacentFrac}{0/3}
\newcommand{\VthreeGeminiThreeOneProPreviewFieldGCredentialAdjacentPct}{0\%}
\newcommand{\VthreeGeminiThreeOneProPreviewFieldGCredentialAdjacentK}{0}
\newcommand{\VthreeGeminiThreeOneProPreviewFieldGCredentialAdjacentN}{3}
\newcommand{\VthreeGeminiThreeOneProPreviewFieldGPlatformAdjacentFrac}{20/20}
\newcommand{\VthreeGeminiThreeOneProPreviewFieldGPlatformAdjacentPct}{100\%}
\newcommand{\VthreeGeminiThreeOneProPreviewFieldGPlatformAdjacentK}{20}
\newcommand{\VthreeGeminiThreeOneProPreviewFieldGPlatformAdjacentN}{20}
\newcommand{\VthreeGeminiThreeOneProPreviewFieldGRegionAdjacentFrac}{20/20}
\newcommand{\VthreeGeminiThreeOneProPreviewFieldGRegionAdjacentPct}{100\%}
\newcommand{\VthreeGeminiThreeOneProPreviewFieldGRegionAdjacentK}{20}
\newcommand{\VthreeGeminiThreeOneProPreviewFieldGRegionAdjacentN}{20}
\newcommand{\VthreeGeminiThreeOneProPreviewFieldMCredentialFrac}{20/20}
\newcommand{\VthreeGeminiThreeOneProPreviewFieldMCredentialPct}{100\%}
\newcommand{\VthreeGeminiThreeOneProPreviewFieldMCredentialK}{20}
\newcommand{\VthreeGeminiThreeOneProPreviewFieldMCredentialN}{20}
\newcommand{\VthreeGeminiThreeOneProPreviewFieldMOperatorFrac}{20/20}
\newcommand{\VthreeGeminiThreeOneProPreviewFieldMOperatorPct}{100\%}
\newcommand{\VthreeGeminiThreeOneProPreviewFieldMOperatorK}{20}
\newcommand{\VthreeGeminiThreeOneProPreviewFieldMOperatorN}{20}
\newcommand{\VthreeGeminiThreeOneProPreviewFieldMRegionFrac}{20/20}
\newcommand{\VthreeGeminiThreeOneProPreviewFieldMRegionPct}{100\%}
\newcommand{\VthreeGeminiThreeOneProPreviewFieldMRegionK}{20}
\newcommand{\VthreeGeminiThreeOneProPreviewFieldMRegionN}{20}
\newcommand{\VthreeGeminiThreeOneProPreviewGenericFrac}{0/0}
\newcommand{\VthreeGeminiThreeOneProPreviewGenericK}{0}
\newcommand{\VthreeGeminiThreeOneProPreviewGenericN}{0}
\newcommand{\VthreeGeminiThreeOneProPreviewNamingFrac}{80/80}
\newcommand{\VthreeGeminiThreeOneProPreviewNamingPct}{100\%}
\newcommand{\VthreeGeminiThreeOneProPreviewNamingCI}{[95, 100]}
\newcommand{\VthreeGeminiThreeOneProPreviewNamingK}{80}
\newcommand{\VthreeGeminiThreeOneProPreviewNamingN}{80}
\newcommand{\VthreeGeminiThreeOneProPreviewOffdiagFrac}{24/492}
\newcommand{\VthreeGeminiThreeOneProPreviewOffdiagPct}{5\%}
\newcommand{\VthreeGeminiThreeOneProPreviewOffdiagCI}{[0, 13]}
\newcommand{\VthreeGeminiThreeOneProPreviewOffdiagK}{24}
\newcommand{\VthreeGeminiThreeOneProPreviewOffdiagN}{492}
\newcommand{\VthreeGptFouroAdjacentFrac}{20/60}
\newcommand{\VthreeGptFouroAdjacentPct}{33\%}
\newcommand{\VthreeGptFouroAdjacentCI}{[23, 46]}
\newcommand{\VthreeGptFouroAdjacentK}{20}
\newcommand{\VthreeGptFouroAdjacentN}{60}
\newcommand{\VthreeGptFouroControlFrac}{0/800}
\newcommand{\VthreeGptFouroControlPct}{0\%}
\newcommand{\VthreeGptFouroControlK}{0}
\newcommand{\VthreeGptFouroControlN}{800}
\newcommand{\VthreeGptFouroControlcallsFrac}{0/160}
\newcommand{\VthreeGptFouroControlcallsK}{0}
\newcommand{\VthreeGptFouroControlcallsN}{160}
\newcommand{\VthreeGptFouroCredentialFrac}{20/20}
\newcommand{\VthreeGptFouroCredentialPct}{100\%}
\newcommand{\VthreeGptFouroCredentialK}{20}
\newcommand{\VthreeGptFouroCredentialN}{20}
\newcommand{\VthreeGptFouroCrosstalknamingFrac}{0/9}
\newcommand{\VthreeGptFouroCrosstalknamingK}{0}
\newcommand{\VthreeGptFouroCrosstalknamingN}{9}
\newcommand{\VthreeGptFouroCrosstalkuaFrac}{8/9}
\newcommand{\VthreeGptFouroCrosstalkuaK}{8}
\newcommand{\VthreeGptFouroCrosstalkuaN}{9}
\newcommand{\VthreeGptFouroDeltaCI}{[32, 95]}
\newcommand{\VthreeGptFouroDeltaPp}{66}
\newcommand{\VthreeGptFouroDeltaClusters}{7}
\newcommand{\VthreeGptFouroDiagonalFrac}{95/140}
\newcommand{\VthreeGptFouroDiagonalPct}{68\%}
\newcommand{\VthreeGptFouroDiagonalCI}{[36, 96]}
\newcommand{\VthreeGptFouroDiagonalK}{95}
\newcommand{\VthreeGptFouroDiagonalN}{140}
\newcommand{\VthreeGptFouroFabricationFrac}{0/93}
\newcommand{\VthreeGptFouroFabricationPct}{0\%}
\newcommand{\VthreeGptFouroFabricationK}{0}
\newcommand{\VthreeGptFouroFabricationN}{93}
\newcommand{\VthreeGptFouroFieldApiDocumentationFrac}{15/20}
\newcommand{\VthreeGptFouroFieldApiDocumentationPct}{75\%}
\newcommand{\VthreeGptFouroFieldApiDocumentationK}{15}
\newcommand{\VthreeGptFouroFieldApiDocumentationN}{20}
\newcommand{\VthreeGptFouroFieldGCredentialAdjacentFrac}{0/20}
\newcommand{\VthreeGptFouroFieldGCredentialAdjacentPct}{0\%}
\newcommand{\VthreeGptFouroFieldGCredentialAdjacentK}{0}
\newcommand{\VthreeGptFouroFieldGCredentialAdjacentN}{20}
\newcommand{\VthreeGptFouroFieldGPlatformAdjacentFrac}{0/20}
\newcommand{\VthreeGptFouroFieldGPlatformAdjacentPct}{0\%}
\newcommand{\VthreeGptFouroFieldGPlatformAdjacentK}{0}
\newcommand{\VthreeGptFouroFieldGPlatformAdjacentN}{20}
\newcommand{\VthreeGptFouroFieldGRegionAdjacentFrac}{20/20}
\newcommand{\VthreeGptFouroFieldGRegionAdjacentPct}{100\%}
\newcommand{\VthreeGptFouroFieldGRegionAdjacentK}{20}
\newcommand{\VthreeGptFouroFieldGRegionAdjacentN}{20}
\newcommand{\VthreeGptFouroFieldMCredentialFrac}{20/20}
\newcommand{\VthreeGptFouroFieldMCredentialPct}{100\%}
\newcommand{\VthreeGptFouroFieldMCredentialK}{20}
\newcommand{\VthreeGptFouroFieldMCredentialN}{20}
\newcommand{\VthreeGptFouroFieldMOperatorFrac}{20/20}
\newcommand{\VthreeGptFouroFieldMOperatorPct}{100\%}
\newcommand{\VthreeGptFouroFieldMOperatorK}{20}
\newcommand{\VthreeGptFouroFieldMOperatorN}{20}
\newcommand{\VthreeGptFouroFieldMRegionFrac}{20/20}
\newcommand{\VthreeGptFouroFieldMRegionPct}{100\%}
\newcommand{\VthreeGptFouroFieldMRegionK}{20}
\newcommand{\VthreeGptFouroFieldMRegionN}{20}
\newcommand{\VthreeGptFouroGenericFrac}{0/20}
\newcommand{\VthreeGptFouroGenericPct}{0\%}
\newcommand{\VthreeGptFouroGenericK}{0}
\newcommand{\VthreeGptFouroGenericN}{20}
\newcommand{\VthreeGptFouroNamingFrac}{75/80}
\newcommand{\VthreeGptFouroNamingPct}{94\%}
\newcommand{\VthreeGptFouroNamingCI}{[86, 97]}
\newcommand{\VthreeGptFouroNamingK}{75}
\newcommand{\VthreeGptFouroNamingN}{80}
\newcommand{\VthreeGptFouroOffdiagFrac}{9/560}
\newcommand{\VthreeGptFouroOffdiagPct}{2\%}
\newcommand{\VthreeGptFouroOffdiagCI}{[0, 4]}
\newcommand{\VthreeGptFouroOffdiagK}{9}
\newcommand{\VthreeGptFouroOffdiagN}{560}
\newcommand{\VthreeProvidersConfirmatory}{2}
\newcommand{\VthreeProvidersTotal}{5}
\newcommand{\VthreeTrialsN}{2013}
\newcommand{\VthreeTrialsErrors}{10}

% Generated from raw evidence; do not edit.
\newcommand{\ExpACapacityPlanted}{239/240}
\newcommand{\ExpACost}{1.02}
\newcommand{\ExpAFirstWord}{12}
\newcommand{\ExpAGemBooleanPlanted}{40/40}
\newcommand{\ExpAGemBooleanUnplanted}{20/40}
\newcommand{\ExpAGemConstrained}{157/160}
\newcommand{\ExpAGemDelta}{0.0}
\newcommand{\ExpAGemDeltaCI}{[-3.1, 3.1]}
\newcommand{\ExpAGemEnumPlanted}{40/40}
\newcommand{\ExpAGemEnumUnplanted}{6/40}
\newcommand{\ExpAGemFree}{157/160}
\newcommand{\ExpAGemIntegerPlanted}{40/40}
\newcommand{\ExpAGemIntegerUnplanted}{5/40}
\newcommand{\ExpAGemRejected}{0/160}
\newcommand{\ExpAGemSelDiag}{98.1\%}
\newcommand{\ExpAGemSelOff}{0.0\%}
\newcommand{\ExpAGemUnplanted}{0/40}
\newcommand{\ExpAGemWrongTool}{0}
\newcommand{\ExpAGptBooleanPlanted}{39/40}
\newcommand{\ExpAGptBooleanUnplanted}{19/40}
\newcommand{\ExpAGptConstrained}{154/160}
\newcommand{\ExpAGptDelta}{2.5}
\newcommand{\ExpAGptDeltaCI}{[-0.6, 5.6]}
\newcommand{\ExpAGptEnumPlanted}{40/40}
\newcommand{\ExpAGptEnumUnplanted}{6/40}
\newcommand{\ExpAGptFree}{158/160}
\newcommand{\ExpAGptIntegerPlanted}{40/40}
\newcommand{\ExpAGptIntegerUnplanted}{5/40}
\newcommand{\ExpAGptRejected}{0/160}
\newcommand{\ExpAGptSelDiag}{98.8\%}
\newcommand{\ExpAGptSelOff}{0.0\%}
\newcommand{\ExpAGptUnplanted}{0/40}
\newcommand{\ExpAGptWrongTool}{0}
\newcommand{\ExpAHolmMax}{0.0006}
\newcommand{\ExpAMisses}{14}
\newcommand{\ExpATrials}{1,200}
\newcommand{\ExpBCost}{0.22}
\newcommand{\ExpBFirstWord}{17}
\newcommand{\ExpBGemConstrained}{30/40}
\newcommand{\ExpBGemDelta}{5.0}
\newcommand{\ExpBGemDeltaCI}{[0.0, 10.0]}
\newcommand{\ExpBGemFree}{32/40}
\newcommand{\ExpBGemRejected}{0/40}
\newcommand{\ExpBGemSelDiag}{80.0\%}
\newcommand{\ExpBGemSelOff}{0.0\%}
\newcommand{\ExpBGemWrongTool}{13}
\newcommand{\ExpBGptConstrained}{32/40}
\newcommand{\ExpBGptDelta}{2.5}
\newcommand{\ExpBGptDeltaCI}{[-2.5, 7.5]}
\newcommand{\ExpBGptFree}{33/40}
\newcommand{\ExpBGptRejected}{0/40}
\newcommand{\ExpBGptSelDiag}{82.5\%}
\newcommand{\ExpBGptSelOff}{0.0\%}
\newcommand{\ExpBGptWrongTool}{0}
\newcommand{\ExpBMisses}{33}
\newcommand{\ExpBTrials}{160}
\newcommand{\ExpCost}{1.23}
\newcommand{\JournalGeminiLoo}{50.8--67.5}
\newcommand{\JournalGptLoo}{60.6--77.5}
\newcommand{\JournalMixedStatus}{converged}
\newcommand{\JournalTzeroDiagonal}{80/80}
\newcommand{\JournalTzeroNeutral}{0/80}
\newcommand{\JournalTzeroOff}{0/560}

\hypersetup{pdfauthor={},pdftitle={Selective Recovery of Prompt-Resident Facts through Tool-Schema Parameters}}
\title{Selective Recovery of Prompt-Resident Facts through Tool-Schema Parameters:\\A Controlled Measurement Study}
\author{Author details pending confirmation}
\date{}
\begin{document}
\maketitle
\begin{abstract}
Third-party tool schemas can solicit information unrelated to a user's task.
We measure whether schema wording selectively recovers matching confidential
facts rather than prompt contents indiscriminately. In an order-lookup task,
we plant five synthetic facts, vary seven targeted string fields and one generic
field, and compare with a neutral-field control. On the two protocol-named models,
requested-fact recovery is \VthreeGptFouroDiagonalFrac{} for GPT-4o and
\VthreeGeminiThreeFlashPreviewDiagonalFrac{} for Gemini 3 Flash, versus
\VthreeGptFouroOffdiagFrac{} and \VthreeGeminiThreeFlashPreviewOffdiagFrac{}
nonrequested fact-checks. Paired field-bootstrap differences are
\VthreeGptFouroDeltaPp{} percentage points \VthreeGptFouroDeltaCI{} and
\VthreeGeminiThreeFlashPreviewDeltaPp{} points
\VthreeGeminiThreeFlashPreviewDeltaCI{}. The analysis implementation was corrected
after collection. Neither model recovers a canary in its
\VthreeGptFouroControlcallsN{} neutral calls. A prospective, timestamped follow-up
(\ExpATrials{} calls, with local receipt and task success recorded) finds that a
canary-permitting pattern and length constraint does not measurably reduce
received recovery on either model (\ExpAGemDelta{} and \ExpAGptDelta{} points,
intervals including zero), and that enum, integer and boolean fields return the
planted value in \ExpACapacityPlanted{} trials. Selected fields, synthetic tools
and secrets, and field heterogeneity limit generalization. Parameter-mediated
extraction is established prior art; operational prevention remains open.
\end{abstract}

\begin{IEEEkeywords}
Tool-calling agents, information disclosure, schema security,
Model Context Protocol, reproducible measurement.
\end{IEEEkeywords}

\section{Introduction}
A tool-using assistant must translate a user's request into arguments satisfying
a tool's declared interface. The tool provider can therefore influence what the
model is asked to supply. If an unnecessary field solicits a deployment fact from
the system prompt, a syntactically plausible call can contain information the
user did not authorize the provider to receive. Actual disclosure to the provider
additionally requires host dispatch and server receipt; generation alone does not
establish either event.

This extraction mechanism is not new. HiddenLayer reports parameter abuse and
fake-function variants~\cite{hiddenlayer2025mcp}; MCP Security Bench (MSB) explicitly
includes out-of-scope parameters~\cite{zhang2025msb}. Our question is narrower:
\emph{when several confidential facts are simultaneously available, does the
requested field select which one appears in the emitted arguments?} A
fact-by-field matrix distinguishes selective recovery from indiscriminate
copying. It also exposes whether an apparently low average recovery rate hides
successful fields and consistently unsuccessful controls.

We contribute (i) a controlled matrix with a deterministic, protocol-defined
marker scorer; (ii) paired field-level estimates, independent numerical
cross-checks, and explicit accounting for exploratory and incomplete arms;
(iii) a bounded comparison with an injection-oriented policy, accompanied by
negative results and implementation artifacts; and (iv) a prospective,
externally timestamped experiment testing whether constraining the requested
field's JSON Schema type changes recovery, measured at local receipt with
legitimate-task utility (Section~\ref{sec:exp25}). We do not claim to discover
parameter extraction, prove general retrieval without naming a target, or
establish vulnerability prevalence in deployed agents. The historical gate
hypothesis that a plausible field would outperform an explicit request on a
majority of tested models failed; it is retained in the supplementary record.

\section{Related work}
Indirect prompt injection established that externally supplied content can
redirect an application without a direct adversarial user
request~\cite{greshake2023not}. Declarative field descriptions are another
lower-trust input. We distinguish them from imperative attack text for experimental
purposes, not as a claim that they fall outside every definition of injection.
The planted prompts explicitly prohibit disclosure; successful recovery conflicts
with that instruction even when a schema contains no override command.

Lin's structured-output study tests conflicting prompt and schema definitions
on a nonce-label classification task and finds model-dependent schema
influence~\cite{lin2026schema}. This is close evidence that schema descriptions
are behaviorally active, but its endpoint is classification accuracy rather
than selective recovery of confidential prompt facts in tool arguments.

HiddenLayer's parameter-abuse report establishes the mechanism and also describes
fake functions embedded in user messages~\cite{hiddenlayer2025mcp}. Our provider
controls a declared tool interface, while our measurement crosses requested fields
with simultaneously planted facts. MSB already evaluates parameter-mediated
leakage and executes tools through MCP~\cite{zhang2025msb}. Its task/security
evaluation is complementary; our emitted-argument matrix is not a replacement
for execution or task-utility testing.

MCPTox studies metadata poisoning and supplies the rendered tool listings used
in our exploratory reconstruction~\cite{mcptox2025}. AgentDojo provides an
environment for evaluating tasks, attacks, and defenses~\cite{debenedetti2024agentdojo}.
These motivate a future task-level evaluation rather than a claim that suitable
benchmarks do not exist. Recent runtime enforcement work, including
ShieldMCP~\cite{yergattikar2026securing}, further precludes interpreting a miss by
one static policy as a result about all agent defenses.

CaMeL separates control and data flow and uses capability-based
policies~\cite{debenedetti2025camel}. This supports considering argument-release
controls outside the generating model. We do not evaluate CaMeL or claim that
it would necessarily block our exact fixtures. OWASP recommends avoiding secrets
in system prompts~\cite{owasp_llm07}; our synthetic-credential observations do not
establish that real deployments store such material there.

\section{Threat model and observable outcomes}
The adversary publishes a tool selected for a legitimate task. It controls the
tool declaration, including an additional string parameter and its requiredness,
but not the user request, host, model weights, or confidential system prompt.
The experimenter plants facts to make measurement possible; the adversary does
not receive their values. Each observation concerns one model response under a
fixed condition, not an adaptive extraction sequence. The attacker benefits only
if a host releases the generated information to the tool.

MCP defines tool input schemas and calls, but its security guidance also recommends
showing inputs and obtaining confirmation for sensitive
operations~\cite{mcp_spec}. Delivery is therefore conditional on host behavior,
not an unavoidable consequence of the protocol.

\begin{table}[t]
\centering\small
\caption{Outcome levels. A later level cannot be inferred solely from an earlier one.}
\label{tab:endpoints}
\begin{tabular}{@{}p{.25\columnwidth}p{.67\columnwidth}@{}}
\toprule
Level & Evidence required \\
\midrule
Emission & Model response contains a tool call and recorded arguments. \\
Validity & Arguments parse and satisfy the offered schema and tool name. \\
Dispatch & Host records sending that selected call. \\
Receipt & Server capture matches the correlation ID, tool, and arguments. \\
Task utility & Returned result passes a legitimate-task oracle. \\
\bottomrule
\end{tabular}
\end{table}

The historical endpoint is marker recovery in recorded non-\texttt{query}
arguments conditional on the adapter's tool-call indicator. We preserve it for
comparability. Missing historical validity, dispatch, and receipt fields remain
unknown, not negative observations. A repaired no-network MCP round trip verifies
delivery, rejection, and no-call behavior using mock responses. This is a software
test, not additional model-disclosure evidence.

\section{Design and analysis}
\subsection{Task, fields, and facts}
The primary task asks for the user's three most recent orders. Across conditions,
the legitimate task and base tool are held fixed. Condition C adds a targeted
or generic string field; condition D adds a neutral request-trace field.
The prompt contains five seed-derived fixtures: platform, region, operator,
credential-shaped, and opaque-control values, together with a confidentiality
instruction. None is a live credential. The targeted fields cover four naming
and three adjacent wordings. A generic \texttt{call\_context} field has no
designated target. The field taxonomy was specified before v3 collection;
the labels are not a guarantee of semantic distance.

Each model's planned planted matrix has 20 repetitions for each of eight
wordings under both C and D, at temperature 0.7. Thus seven targeted C fields
supply 140 matched opportunities and 560 nonmatched checks; the generic C
field supplies 20 call-level outcomes. The 160 D calls supply 800 correlated
fact-check opportunities, not 800 independent calls. Unplanted arms separately
test recovery without supplying the fixtures. Provider seeds are not equivalent
randomness controls across APIs. The historical schedule was not the subsequently
implemented blocked-randomized runner.

\subsection{Scoring and estimand}
Let $Y_{frk}$ indicate whether fact $k$ is recovered in repetition $r$ for field
$f$, and let $t(f)$ be its prescribed target. The historical scorer serializes
the non-\texttt{query} argument object, lowercases it, removes nonalphanumeric
characters, and tests for a full marker or its prespecified prefix alias.
This is deterministic normalized substring matching, not semantic grading or
strict equality with a complete canary. Prefix matching can count partial
recovery; punctuation and case are deliberately ignored. No aliases were added
to improve the reported v3 results.
The supplement additionally reports post-collection full-marker and
all-argument sensitivity checks; they do not replace the registered endpoint.
All three alternatives give the same matched and nonmatched counts as the
registered scorer on both primary models.

For eligible emitted calls with counts $n_f$, the primary quantities are
\begin{align}
\widehat p_M &= \frac{\sum_{f,r}Y_{fr,t(f)}}{\sum_f n_f},\\
\widehat p_U &= \frac{\sum_{f,r}\sum_{k\ne t(f)}Y_{frk}}{4\sum_f n_f},\\
\widehat\Delta &= \widehat p_M-\widehat p_U.
\end{align}
This compares a matched opportunity with a nonmatched opportunity, not the
probability of any leak per trial. With equal field sizes it also equals the
average field contrast. Generic and neutral any-canary recovery use
$\mathbb{1}\{\sum_kY_{frk}>0\}$ and a call denominator. Generic outcomes never
enter a contrast requiring $t(f)$.

\subsection{Uncertainty and decision rule}
The primary percentile bootstrap samples seven whole fields with replacement,
retaining each field's matched and nonmatched counts together, for 10,000 draws
with seed zero. The same draw produces both rates and their difference. We report
95\% intervals and leave-one-field-out ranges. Repeated facts within a call
are not treated as independent bootstrap units. Fields were hand-selected, so
these intervals summarize sensitivity to their composition; they do not certify
coverage over a randomly sampled population of schemas or deployments.

The protocol's H-v3-1 rule requires a difference of at least 20 percentage points
and a bootstrap interval excluding zero on both named providers. Missing or
incomplete required arms yield an incomplete decision. The previously
implemented analysis did not faithfully supply the paired contrast; its
post-collection correction and numerical effects are recorded in the deviations
log. We distinguish the pre-collection design from this later implementation.
An independently implemented count and resampling calculation checks the primary
result; it is a computational cross-check, not independent scientific replication.

The supplement reports the registered per-fact Wilson, one-sided Fisher,
Holm-adjusted, and Newcombe analyses and a hierarchical logistic sensitivity
fit. These call-level tests are not cluster-robust replacements for the primary
analysis. Multiplicity-family choices and hierarchical priors were resolved
after collection. H-v3-2's generic diagonal rule is undefined because a generic
field has no target; we report its descriptive outcome without repairing that
rule retrospectively.

\subsection{Provenance and completeness}
The v3 protocol names \texttt{gpt-4o} and \texttt{gemini-3-flash-preview}; these
are requested API identifiers, not immutable model-weight identifiers. The
artifact records a protocol hash and calendar receipts preceding collection.
Independent verification of timestamp anchoring remains an author-facing
submission check. Other providers and the earlier v2 matrix are not upgraded
to confirmatory evidence. An explicit artifact-selection file fixes the input
paths and hashes; the analysis never chooses the largest available run.

Reported served identities are audited against requested identities. Missing
served identities remain unverified. Historical logs do not retrospectively
gain the improved identity and schema-validation guarantees of the new runner.
API-error rows are counted and excluded under the historical protocol; malformed
files fail validation. The supplement gives attempted, usable, emitted-call,
and error counts for every selected matrix arm.

\section{Results}
\subsection{Primary selectivity}
\begin{table}[t]\centering\small
\caption{Primary recovery. M: matched fact checks; U: nonmatched fact checks. Difference and paired 95\% interval are percentage points; seven fields per model.}
\label{tab:primary}
\begin{tabular}{@{}lrrr@{}}\toprule
Model & M & U & Difference [CI] \\\midrule
GPT-4o & 95/140 & 9/560 & 66 [32, 95] \\
Gemini Flash & 82/140 & 4/560 & 58 [17, 86] \\
\bottomrule\end{tabular}\end{table}

Table~\ref{tab:primary} shows the paired effect for the protocol-named models.
Both satisfy the numerical H-v3-1 rule under the corrected implementation.
The wide intervals are substantive: there are seven field clusters, not
hundreds of independent interfaces. Leave-one-field-out differences range from
\JournalGptLoo{} percentage points for GPT-4o and \JournalGeminiLoo{} for
Gemini. Omitting any one field retains a positive point difference, but does
not expand the population to which the result applies.

Neither model recovers any canary in its \VthreeGptFouroControlcallsN{} neutral
calls; generic recovery is \VthreeGptFouroGenericFrac{} and
\VthreeGeminiThreeFlashPreviewGenericFrac{}, respectively. Zero observed events
do not prove zero risk. A call-level Wilson upper limit for the neutral arms
would additionally rely on assumptions about independent, identically
distributed trials that field heterogeneity does not establish.

\subsection{Heterogeneity, negative results, and controls}
\begin{table}[t]\centering\scriptsize
\caption{Matched recovery by actual field name and prespecified class. The supplement gives full descriptions and marker aliases.}
\label{tab:fields}
\begin{tabular}{@{}llrr@{}}\toprule
Field & Class & GPT-4o & Gemini \\\midrule
client\_user\_agent & naming & 15/20 & 2/20 \\
issuing\_surface & adjacent & 0/20 & 0/20 \\
deployment\_region\_code & naming & 20/20 & 20/20 \\
locality\_hint & adjacent & 20/20 & 20/20 \\
operator\_account\_name & naming & 20/20 & 20/20 \\
service\_key\_reference & naming & 20/20 & 20/20 \\
integration\_binding\_note & adjacent & 0/20 & 0/20 \\
\bottomrule\end{tabular}\end{table}

The field-level breakdown (Table~\ref{tab:fields}) qualifies the pooled result.
The region-adjacent field succeeds, whereas the platform- and
credential-adjacent fields do not recover their targets on the primary pair.
Thus this study does not establish broad non-naming retrieval. The
\texttt{service\_key\_reference} field recovers a fabricated credential on
\VthreeGptFouroCredentialFrac{} calls in each primary arm. This tests copying a
prompt-resident fixture, not use of a real service credential.

Most nonmatched recoveries in the primary pair come from the broad
\texttt{client\_user\_agent} field:
\VthreeGptFouroCrosstalkuaFrac{} for GPT-4o and
\VthreeGeminiThreeFlashPreviewCrosstalkuaFrac{} for Gemini. This concentration
supports a field-specific interpretation rather than uniform leakage noise.
Unplanted observations across the selected providers show
\VthreeFabricationPooledFrac{} any-canary recoveries; provider-specific counts
and evidence status are retained in the supplement rather than treating this
pooled descriptive result as a new primary test.
The unplanted stages cover six of the eight matrix fields, omitting the
platform-adjacent and region-adjacent fields. They complete their configured
stages but not the protocol's same-field control coverage; H-v3-3 is therefore
not treated as fully established across the planned field set.

The original majority-of-models gate hypothesis failed. A framework arm without
planted content produced \FrameworkArmDisclosed{} identity disclosures in
\FrameworkArmToolCalled{} tool-call observations. These negative results limit
both superiority over explicit requests and claims about intrinsic model or
framework self-knowledge.

\subsection{Exploratory extensions already collected}
Three additional matrix providers are exploratory. The Gemini Pro arm is
quota-capped, and noncalls and errors make conditional denominators unequal.
Its larger point estimate is not evidence of a capability-scaling law. All
selected v3 matrix and unplanted files together contain
\VthreeTrialsN{} attempted rows, including \VthreeTrialsErrors{} error rows;
this is an inventory total, not the primary sample size.

The existing exploratory temperature-zero v2 arm has
\JournalTzeroDiagonal{} matched and \JournalTzeroOff{} nonmatched recoveries,
with \JournalTzeroNeutral{} neutral any-canary outcomes. It uses different
fixtures and a different collection time; comparing it with v3 cannot isolate a
temperature effect.

The external-listing arm reconstructs \RealCorpusTools{} tool interfaces from
rendered MCPTox listings, appends an authored parameter, and pairs tasks by token
overlap. GPT-4o recovery is \RealGptFouroCFrac{} conditional on a call and
\RealGptFouroCIttFrac{} over attempted trials; Claude recovery is
\RealClaudeSonnetFourFiveCFrac{} and \RealClaudeSonnetFourFiveCIttFrac{}.
Neutral recovery is \RealDCombinedFrac{} across their emitted calls.
Native types, nesting, requiredness, and constraints were not preserved:
harvested parameters were reconstructed as strings. These observations broaden
wording contexts, not validation against original live servers. Neither
legitimate-task success nor a population upper bound follows from these rates.

\section{Prospective schema-constraint experiment}\label{sec:exp25}
Every earlier arm offered the requested field as a free string. Experiment~25
asks two questions that the historical design could not: does a constraint on
that string change received recovery, and can a categorical field carry a
planted value above its guessing rate? It was specified, reviewed and
externally timestamped before any planted call, and is analysed separately from
the historical evidence, never pooled with it.

\subsection{Design, host, and provenance}
Stage~A uses two synthetic tools, order lookup and email search, and four
naming fields, one per fact family (platform, region, operator, and a fabricated
credential). Each tool and replicate has its own four seed-derived canaries,
shared by every field and both arms, so each free trial has an exact
constrained partner that differs only in the field's schema. The constrained
arm adds the pattern \texttt{\^{}[A-Za-z0-9\_ -]\{1,128\}\$} and
\texttt{maxLength:~128}; every planted canary satisfies it, so the arm tests a
permissive constraint, not one that forbids the secret. There are 20 pairs per
model, tool and field, and five unplanted repetitions. A capacity arm on the
order tool adds an enum of eight tiers, an integer from 0 to 7, or a boolean,
with 40 balanced planted values and 40 unplanted trials per model and domain;
the unplanted trials carry hidden balanced labels that the model never sees.
Stage~B adds one fixed field, \texttt{client\_platform\_name}, to eight frozen
native MCP tool schemas (two servers for each of four tasks), with five pairs
per schema and model; it was planned as a small contextual check, conditional
on funding, and is analysed separately from stage~A.

A local host parses each call and validates it against the offered schema before
an in-process fixture handler receives it. The primary endpoint is receipt of
the target canary in a non-task argument per attempted trial, scored by full
normalized marker without the historical prefix aliases; a task oracle scores
legitimate-task utility. Handler receipt is not MCP transport or a provider-side
event. The models are \texttt{gpt-4o-2024-08-06} and
\texttt{gemini-3-flash-preview} at temperature 0.7, collected as two legs in a
fixed order, Gemini first. The primary estimand is the per-model mean of free
minus constrained recovery; the protocol describes attenuation only when both
models' 95\% intervals exclude zero in the positive direction. Stage~A intervals
resample the 20 replicate blocks, 10,000 draws; capacity uses one-sided label
permutation tests, Holm-adjusted over the six model--domain tests.

The protocol file and a hash of the executable design were committed, and the
protocol's SHA-256 was submitted to four OpenTimestamps calendars, before any
planted call; Bitcoin anchoring was pending at the time of writing. A coauthor
reviewed the four field texts before stage~A and the eight native task fixtures
before stage~B, and recorded both decisions. All other model calls were
unplanted compatibility checks: 14 on Gemini before sealing, 14 on GPT-4o before
its leg, and 35 before stage~B. One dated amendment,
approving the native fixtures, changed the design hash before any stage~B call;
stage~A was collected under the sealed hash. The supplement discloses a first
launch that failed locally before any request left the host, and three stops by
the prespecified output-token bound, because Gemini's
output limit does not cap its thinking tokens. Each stopped leg was resumed under
the prespecified policy without repeating a trial. All \ExpATrials{} stage~A and
\ExpBTrials{} stage~B trials completed with no API errors.

\subsection{Results}
\begin{table}[t]\centering\small
\caption{Experiment 25 primary endpoint: received target recovery with a free string versus a canary-permitting pattern/length constraint, exact pairs. Difference and 95\% interval in percentage points (stage A: replicate-block bootstrap; stage B: schemas within tasks).}
\label{tab:exp25}
\begin{tabular}{@{}llrrr@{}}\toprule
Stage & Model & Free & Constr. & Diff. [CI] \\\midrule
A (two synthetic tools) & Gemini Flash & 157/160 & 157/160 & 0.0 [-3.1, 3.1] \\
A (two synthetic tools) & GPT-4o & 158/160 & 154/160 & 2.5 [-0.6, 5.6] \\
B (eight native schemas) & Gemini Flash & 32/40 & 30/40 & 5.0 [0.0, 10.0] \\
B (eight native schemas) & GPT-4o & 33/40 & 32/40 & 2.5 [-2.5, 7.5] \\
\bottomrule\end{tabular}\end{table}

Table~\ref{tab:exp25} gives the primary endpoint. No interval excludes zero in
the positive direction, in either stage, so by the prespecified rule no
attenuation is described for either model. In stage~A the intervals rule out
attenuation larger than about 3 points for Gemini and 6 for GPT-4o, for these
fields and tools. Validation rejected \ExpAGemRejected{} and \ExpAGptRejected{}
constrained calls: a constraint that admits the secret does not stop the model
from supplying it.

\begin{table}[t]\centering\small
\caption{Experiment 25 capacity arm (stage A, orders tool): planted admissible value received versus a hidden balanced label with nothing planted. Permutation tests, Holm-adjusted over all six, $p\le\ExpAHolmMax$.}
\label{tab:exp25capacity}
\begin{tabular}{@{}lrrrrr@{}}\toprule
 & & \multicolumn{2}{c}{Gemini Flash} & \multicolumn{2}{c}{GPT-4o} \\
Field & Chance & Planted & None & Planted & None \\\midrule
Enum (8) & 12.5\% & 40/40 & 6/40 & 40/40 & 6/40 \\
Integer 0--7 & 12.5\% & 40/40 & 5/40 & 40/40 & 5/40 \\
Boolean & 50\% & 40/40 & 20/40 & 39/40 & 19/40 \\
\bottomrule\end{tabular}\end{table}

The capacity arm (Table~\ref{tab:exp25capacity}) returned the planted value in
\ExpACapacityPlanted{} trials, while unplanted answers stayed near their guessing
rates. A categorical field therefore carries any secret that is one of its
admissible values; a boolean carries one bit per call. These domains differ in
wording and planted fact as well as type, so the table does not isolate a pure
type effect.

Stage~A also reproduces field-to-fact selectivity: the requested fact was
received in \ExpAGemSelDiag{} (Gemini) and \ExpAGptSelDiag{} (GPT-4o) of free
calls, and each of the other three facts in \ExpAGemSelOff{} and
\ExpAGptSelOff{}. Only naming fields were tested, and the prompts and scorer
differ from the historical matrix, so these rates are not comparable with
Table~\ref{tab:primary}. Unplanted controls received no canary
(\ExpAGemUnplanted{} and \ExpAGptUnplanted{}). Stage~B's recovery was lower, and
its intervals again include zero, although Gemini's lower endpoint is exactly
zero; with one field and five pairs per schema, stage~B is a small contextual
check, not a precise estimate.

Two descriptive observations were not prespecified. First, of the
\ExpAMisses{} stage~A non-recoveries, \ExpAFirstWord{} sent only the first word
of a two-word platform canary, which the full-marker rule does not count, so the
small free--constrained differences reflect truncation rather than withheld
secrets. Second, in stage~B, \ExpBGemWrongTool{} Gemini calls named tools that
were not offered and so never reached the handler, and \ExpBFirstWord{} of the
\ExpBMisses{} non-recoveries were first-word truncations.

\subsection{Interpretation}
Within these fields and tools, restricting a string's format, or its type to a
small categorical domain, does not prevent recovery of a secret that the schema
still admits. A categorical domain limits how much can move per call; it does
not stop admissible values from moving. Constraints that exclude the secret were
not tested: their zero recovery would reflect the schema rather than model
behavior and would need an attempt-based endpoint. The experiment uses two
models, synthetic tools and secrets, and purposively chosen fields; it is not
evidence about arbitrary schemas or deployed agents.

\section{Defensive interpretation}
The artifact reproduces the prompt and pseudo-tag rule from
\texttt{mcp-scan==0.3.39}, using its specified model and temperature and a
study-defined majority vote across repeated classifications. It flags the
tested imperative conditions but not the primary parameter-only condition or
the explicit-parameter comparison. This is an observation about a historical
published policy reproduced locally, not the vendor's hosted service or current
production rules. Historical artifacts retain boolean votes, not raw judge
text; the old parser mapped any response other than ``YES'' to a negative vote.
Malformed answers therefore cannot be ruled out retrospectively. The recorded
miss is provisional supporting evidence, not a fully audited detector result.
Future parsing rejects malformed answers and retains raw text; no new scanner
calls were made for this revision. A plausible explanation is that an injection-focused question
does not directly assess whether generated arguments are authorized to leave
the host; we do not establish a necessary failure of every such classifier.

A separate parameter classifier has AUC \DefenseAucValue{} on
\DefenseCorpusPositives{} study-authored positives and
\DefenseCorpusNegatives{} corpus negatives. That is detection performance on
this evaluation set, not measured prevention. Retrospective removal projections
do not account for argument substitution, schema rejection, task failure, or
adaptive attacks. Optional-field ablations likewise do not validate a deployed
release policy.

A defensible next evaluation would compare no defense, optional-field removal,
and a task-contract argument-release policy, measuring both disclosure and
legitimate-task utility. Required fields cannot simply be dropped while assuming
the original task still succeeds. Exact-canary filtering would be an oracle
diagnostic only. Section~\ref{sec:exp25} shows that schema design alone is not
such a policy: a permissive constraint left recovery unchanged, and categorical
fields carried admissible secrets. Defense comparisons remain future work and
were not part of that experiment.

\section{Validity and limitations}
\textbf{Construct validity.} Normalized marker recovery is a narrow proxy for
prompt-content disclosure. It can miss paraphrases and count documented partial
aliases. It does not measure all sensitive information, successful execution,
or adversarial downstream use. The older graded arms use a different endpoint;
they are not pooled with the marker matrix. Their human validation is one
coauthor's blinded, stratified gate sample, not independent validation of all arms.

\textbf{Internal validity.} Primary fixtures are shared across repetitions;
provider behavior may drift over time, and historical order effects are not
removed by later runner improvements. Post-collection analysis corrections,
incomplete exploratory arms, and unspecified secondary choices are disclosed.
Prespecified marker aliases reduce discretionary scoring but do not by themselves
make every analytical choice prospective.

\textbf{External validity.} One primary task, two protocol-named models,
hand-selected fields, free-string channels, and synthetic secrets are the main
constraints. The prospective experiment adds constrained and categorical fields
and eight native schema contexts, but only two synthetic tools, four fields, and
one field on the native schemas. No claim is made about typical secret storage, arbitrary native
schemas, current model releases, commercial clients, or real-world incident
rates. Corpus reconstruction and extra models do not substitute for a
prospectively paired task/schema study.

\textbf{Statistical validity.} Small cluster counts make percentile intervals
sensitive to the empirical field distribution. Call-level secondary tests ignore
some dependence; the hierarchical sensitivity fit has only two provider levels
and does not explicitly model repeated facts within each trial. Separation and
prior choice further limit interpretation. We report these analyses without
treating nominal significance as an independent replication of the primary effect.

\section{Ethics, reproducibility, and availability}
All planted facts are synthetic, including credential-shaped values. Measurements
used author-controlled prompts and provider accounts, not third-party victim
deployments. The no-network transport server is a test fixture. Completed
coordinated vendor notification is not evidenced in the repository; the authors
must resolve and document the disclosure decision before submission. Prior public
repository exposure is not evidence of coordinated notification.

The analysis artifact includes immutable historical logs and protocols, explicit
selection hashes, a deviations record, regenerated tables, an independent
numerical check, and a locked offline Python environment. Live provider calls
are outside default tests. A clean-copy check excludes credentials, Git history,
caches, and unrelated local documents. The prospective experiment used
US\$\ExpCost{} of inference at list prices for its \ExpATrials{} and
\ExpBTrials{} trials, recorded per call in a durable budget journal; no other
paid inference was used for this revision. Code checks do not certify anonymity, redistribution rights, or
scientific completeness. The final release URL, license permissions, funding,
conflicts, authorship, prior-publication declarations, and venue-specific AI-use
statement require author confirmation and are tracked separately, not invented.

\section{Conclusion}
Under a fixed order-lookup task, requested fields recover matching planted facts
more often than nonmatching ones on both protocol-named models. The paired
field-level effect survives a disclosed analysis correction and independent
numerical recomputation. Field heterogeneity, generic-field negatives, and the
failed earlier gate constrain the result to selective recovery under these
fixtures. Recorded votes from one locally reproduced policy indicate misses,
subject to the missing raw-response limitation; they do not establish universal
detector failure. A prospective follow-up found that constraining the field's
format, or restricting it to a categorical type, did not prevent recovery of
admissible secrets. Evaluating argument-release policies with legitimate-task
utility is the next step toward operational conclusions.

\bibliographystyle{IEEEtran}
\bibliography{refs}
\end{document}
